\section{Introducción}\label{sec:intro}

Que un salario mínimo beneficie a los trabajadores a quienes se dirige depende de a cuántos de ellos llega realmente. En los países en desarrollo, una gran parte del empleo queda fuera del alcance de la regulación laboral, de modo que un piso salarial legal puede elevar la remuneración de los trabajadores cubiertos, desplazar a algunos de ellos hacia el sector informal no cubierto o funcionar solo como un punto de referencia para salarios que las empresas formales habrían pagado de todos modos. Cuál de estas fuerzas predomina es una pregunta empírica, y su respuesta depende de cuán vinculante es el piso, de qué tan bien se fiscaliza y de con cuánta libertad se mueven los trabajadores entre los márgenes formal e informal. Estudio estas preguntas para el Perú, un país que combina un salario mínimo elevado en relación con la remuneración mediana con una de las tasas de informalidad laboral más altas de América Latina.

Aprovecho el aumento del salario mínimo nacional del Perú (la \emph{Remuneración Mínima Vital}, o RMV) de 930 a 1025 soles, dispuesto por el Decreto Supremo 003-2022-TR y vigente desde el 1 de mayo de 2022. Este aumento nominal de 10.2 por ciento fue el primer cambio de la RMV en cuatro años y se produjo cuando el mercado laboral salía de la pandemia. Un hecho institucional central determina el análisis: el Perú fija un único salario mínimo nacional, sin variación regional ni sectorial en el régimen general. Por lo tanto, no existe variación transversal en el nivel de la política ni una adopción escalonada que aprovechar. Como señala \citet{jimenez_2023} para la reforma de 2016, una identificación creíble debe provenir, en cambio, de una \emph{exposición} diferenciada a una política que cambia para todos al mismo tiempo.

Mi fuente de exposición diferenciada es la educación. En los trimestres previos a la reforma, el salario mensual equivalente mediano de un asalariado privado de baja educación era de 950 soles, casi exactamente el mínimo anterior de 930, y el 57 por ciento de los asalariados de baja educación percibía menos que el \emph{nuevo} piso de 1025 soles en vísperas de la reforma; un aumento de diez por ciento en ese piso resulta, por ello, muy vinculante para este grupo. Los trabajadores de alta educación, en cambio, tienen un salario mediano muy por encima del mínimo y es menos probable que este los alcance, aunque una minoría considerable sigue expuesta a través del trabajo informal o por cuenta propia. Esta brecha de exposición, y no el supuesto de una exposición nula del grupo de comparación, motiva un diseño de diferencias en diferencias que contrasta a los trabajadores de baja educación (aquellos con, a lo sumo, secundaria completa, cerca de dos tercios de la población en edad de trabajar) con los trabajadores de alta educación (aquellos con cualquier nivel de educación superior). El diseño traslada una larga tradición de medición de los efectos del salario mínimo a partir de la variación en la incidencia (\emph{bite}) de la política entre grupos o regiones \citep{card_1992, lee_1999, machin_manning_rahman_2003} a un contexto con una sola reforma nacional y datos de hogares trimestrales y ricos en información.

Utilizo los módulos trimestrales de empleo de uso público de la Encuesta Nacional de Hogares (ENAHO) desde 2021Q1 hasta 2024Q4, que rodean la reforma con cinco trimestres previos y diez trimestres posteriores limpios; el trimestre de transición 2022Q2, en el que la reforma entró en vigor a mitad de trimestre, se excluye de la especificación base y se reincorpora en una prueba de robustez. Los datos son cortes transversales repetidos, y su frecuencia trimestral me permite estimar un estudio de eventos dinámico, contrastar tendencias previas diferenciadas por educación y seguir el efecto durante más de dos años después de la reforma. Examino un conjunto amplio de resultados: salarios reales por hora e ingresos mensuales, la probabilidad de estar empleado y de participar en la fuerza laboral, la incidencia del empleo formal frente al informal, las horas de trabajo y la división del empleo entre trabajo asalariado y trabajo independiente. En todo el análisis pondero con los factores de expansión y agrupo la inferencia por departamento; y como el tratamiento se asigna a un nivel agregado, complemento los errores estándar robustos por conglomerados convencionales con la corrección para muestras pequeñas de \citet{pustejovsky_tipton_2018} y con inferencia por aleatorización.

Surgen tres resultados. Primero, y de manera más robusta, la reforma desplazó la \emph{composición} del empleo de baja educación lejos del trabajo asalariado formal. La probabilidad de que un trabajador de baja educación tenga un empleo formal cae en 4.6 puntos porcentuales, cerca de una quinta parte de la tasa base de formalidad del grupo, de 23 por ciento, y el trabajo independiente aumenta en 3.6 puntos. Las tendencias previas de formalidad de estas estimaciones superan su prueba conjunta, con una pequeña deriva residual que cuantifico y corrijo; los efectos aparecen inmediatamente después de la reforma y persisten, son mayores en los departamentos donde el nuevo piso corta más profundamente la distribución salarial y los confirma un diseño independiente de exposición continua que relaciona los resultados con el índice de Kaitz regional previo a la reforma y admite inferencia por aleatorización basada en el diseño. Las personas enlazadas del Panel ENAHO 2020--2024 revelan el flujo que subyace al stock: los trabajadores formales en su puesto fueron retenidos, mientras que la tasa a la que los trabajadores informales de baja educación ingresaban a empleos formales cayó casi a la mitad, de modo que la reasignación operó mediante la creación de vínculos formales y no mediante su destrucción. Segundo, no hay evidencia de que la reforma elevara la remuneración del grupo, pero no porque el piso sea letra muerta. Dentro del sector formal, donde opera la fiscalización, el pico de ingresos en el mínimo anterior se traslada casi por completo al nuevo, y cae la proporción de asalariados formales que perciben menos que el nuevo piso; entre los asalariados informales esa proporción \emph{aumenta}, y como solo el 13 por ciento de los trabajadores ocupados de baja educación son formales, ambas respuestas se compensan: la proporción agregada por debajo del piso no se mueve, ningún decil de la distribución salarial muestra una ganancia y el efecto estimado sobre el logaritmo del salario real por hora promedio es negativo y no significativo. La conclusión de que no hubo una elevación salarial descansa en la evidencia de cumplimiento y distributiva, cuyos diagnósticos son limpios; la regresión del salario promedio, en sí misma, presenta una prueba de tendencias previas rechazada, impulsada por un único trimestre de inicios de 2021, y soy explícito sobre lo que puede y no puede descartar. Tercero, la tasa de empleo de los trabajadores de baja educación cae en 2.5 puntos porcentuales, pero este margen es el menos robusto: no supera las pruebas placebo en el tiempo, presenta una clara tendencia previa diferenciada asociada a la recuperación de la pandemia y no escala con la exposición regional, por lo que no le doy una interpretación causal. En todo el análisis acoto las conclusiones con el análisis de sensibilidad de \citet{rambachan_roth_2023} y con una batería de pruebas de especificación y placebo. Luego replico el diseño completo en la reforma de enero de 2025, que elevó la RMV a 1130 soles en un período más tranquilo, posterior a la pandemia: allí tampoco aparece respuesta salarial, ni tampoco la reasignación, lo que interpreto como evidencia de que el margen en el que el salario mínimo resulta vinculante depende del estado del mercado laboral.

La contribución es una combinación y no una herramienta nueva, y la planteo sin exagerarla. Hasta donde sé, esta es la primera evaluación de la reforma del salario mínimo peruano de 2022 y el primer estudio de cualquier reforma peruana que utiliza la ENAHO trimestral posterior a la pandemia; la evidencia microeconómica existente se detiene en la reforma de 2016 \citep{jimenez_2023} o en las reformas de inicios de la década de 2000 \citep{jaramillo_lopez_2006, cespedes_2005}. Es también el primero en hacer de la heterogeneidad por educación el objeto central de un estudio de eventos a nivel de trabajador sobre el salario mínimo peruano, mientras que los trabajos previos analizan la heterogeneidad por rango de ingresos o edad, o construyen una dosis de exposición a nivel departamental. Es también, hasta donde sé, el primero en usar el Panel ENAHO enlazado para medir transiciones de formalidad a nivel de trabajador en torno a una reforma del salario mínimo peruano, lo que me permite separar la retención de los trabajadores en su puesto de la entrada de nuevos trabajadores a empleos formales. Y sitúa el margen informal en el centro del análisis, al mostrar que el ajuste opera a través de la composición del empleo y no de la remuneración. No reclamo novedad para la idea de identificación, que es el diseño estándar de exposición o de incidencia, ni para las herramientas econométricas; el valor agregado es la combinación de una reforma nueva y un período posterior a la pandemia relevante para la política pública con un diseño a nivel de trabajador basado en la educación y un tratamiento cuidadoso de la informalidad y de las amenazas a la identificación que tal diseño implica.

El resto del artículo se organiza como sigue. La Sección~\ref{sec:institutional} describe el salario mínimo peruano y la reforma, y la Sección~\ref{sec:literature} sitúa el artículo en la literatura. La Sección~\ref{sec:data} presenta los datos y la construcción de las variables clave, y la Sección~\ref{sec:strategy} expone la estrategia empírica y los supuestos de identificación. La Sección~\ref{sec:results} reporta los resultados principales, la Sección~\ref{sec:robustness} las pruebas de robustez e inferencia y la Sección~\ref{sec:heterogeneity} el análisis de heterogeneidad y distributivo. La Sección~\ref{sec:discussion} discute los mecanismos, la validez externa y las limitaciones, y la Sección~\ref{sec:conclusion} concluye.
